\chapter{代数表示}

\section{代数与模}

代数实际上是在环的基础上添加上模/表示的结构.

\begin{definition}[域上的代数]\index{d@代数!y@域上的代数}
    设$F$是域，$A$是$F$上的一个向量空间且$A$是一个结合环，即
    \[ \forall\, a, b, c \in A, a(bc) = (ab)c \]
    若$A$中作为环的加法和在向量空间中的加法相同且$F-\!\!$空间中的标量乘法与环$A$的乘法可交换：
    \[ \forall\, \lambda \in F, a, b \in A, \lambda(ab) = (\lambda a)b = a(\lambda b) \]
    则称$A$是\colorbf{域$F$上的代数}，简称为$F-\!\!$代数.
\end{definition}

由于$A$既是向量空间又是环，故继承了子代数、商代数、理想等概念，但对于$A$同时具有以上两个空间的性质，因此其子代数也必须是子环和子空间.

\begin{definition}[单代数]\index{d@单代数}
    若$F-\!\!$代数$A$仅有$0$和$A$两个非平凡理想，则称其为\colorbf{单代数}.
\end{definition}

同样可以定义代数之间的运算：

\begin{definition}[代数同态]\index{d@代数同态}
    设$A, B$是$F-\!\!$代数，若$f: A \to B$既是环同态，又是$F-\!\!$线性映射，则称其为\colorbf{$F-\!\!$代数同态}.
\end{definition}

进而，可以类比群表示，定义代数表示.和群表示不同的是，代数中的元素不一定可逆，因此我们不要求值域是$V$上的可逆线性变换，而仅仅是自同态即可.除此之外，$V$上的自同态构成代数，因此需要代数同态.

\begin{definition}[代数的表示]\index{d@代数的表示}
    设$A$是$F-\!\!$代数，$V$是$F-\!\!$向量空间，若存在$F-\!\!$代数同态$\rho : A \to \mathrm{End}_F(V)$，则称$(V, \rho)$是$A$的\colorbf{$F-\!\!$表示}.
\end{definition}

在介绍群表示时，我们曾利用群作用给出等价定义，在代数表示中同样成立.
即定义$\rho(a): V \to V, v \mapsto \rho(a)v \triangleq a \cdot v$，并要求其满足代数同态的要求，如此即可得到模.

\begin{definition}[左$A-\!\!$模]\index{m@模!a@$A$-模}
    设$A$为$F-\!\!$代数，$V$是$F-\!\!$空间，若存在映射$\setjot[-2]\begin{aligned}[t]
        A \times V &\to V \\
        (a,v) &\mapsto a \cdot v,\forall\, a \in A, v \in V
    \end{aligned}$满足：
    \begin{enumerate}[(i)]
        \item $a(v_1 + v_2) = a v_1 + a v_2, \quad\forall\, a \in A, v_1, v_2 \in V$
        \item $(\lambda a) v = a (\lambda v) = \lambda (av), \quad\forall\, a \in A, \lambda \in F, v \in V$
        \item $(a+b) v = av + bv, \quad\forall\, a, b \in A, v \in V$
        \item $(ab)v = a(bv), \quad\forall\, a, b \in A, v \in V$
        \item $1_A \cdot v = v, \quad\forall\, v \in V$
    \end{enumerate}
    则称$V$为\colorbf{左$A-\!\!$模}.
\end{definition}

注意，模根据$A$中元素是左乘还是右乘可分为左模和右模，但二者在根本上是等价的，因此我们仅讨论左模即可.\par
更一般地，有以下给出模的定义：
\begin{definition}[模]\index{m@模!r@$R$-模}
    设$R$是环，$M$是Abel群，若存在映射$\setjot[-2]\begin{aligned}[t]
        R \times M &\to M \\
        (r, m) &\mapsto rm
    \end{aligned}$，满足：
    \begin{enumerate}[(i)]
        \item $r(x+y) = rx+ry, \quad\forall\, r \in R, x, y \in M$
        \item $(r+s)x = rx+sx, \quad\forall\, r, s \in R, x \in M$
        \item $1_R \cdot x = x, \quad\forall\, x \in M$
        \item $(rs)x = r(sx), \quad\forall\, r, s \in R, x \in M$
    \end{enumerate}
    则称$M$是\colorbf{左$R-\!\!$模}.
\end{definition}

\begin{instance}
    向量空间是$F-\!\!$模. \\
    $R=\mathbb{Z}$, $(M, +)$为Abel群，则$M$可视作$\mathbb{Z}-\!\!$模.
    其标量乘法定义为：
    \[ nx = \begin{cases} x + \dots + x, & n > 0 \\ 0, & n = 0 \\ -(-nx), & n < 0 \end{cases} \]
\end{instance}

借助一般模的定义，可将域上的代数推广到环上的代数.

\begin{definition}[环上的代数]\index{d@代数!h@环上的代数}
    设$R$是交换环，$A$是环且满足：
    \begin{enumerate}[label=(\roman*)]
        \item $(A, +)$可视作$R-\!\!$模
        \item $\forall\, a,b \in A, r \in R, r(ab) = (ra)b = a(rb)$
    \end{enumerate}
    则称$A$是 \colorbf{$R-\!\!$代数}.
\end{definition}

\begin{instance}\index{q@群代数}
    考虑群代数$\mathbb{F}G$，其作为$\mathbb{F}-\!\!$线性空间以群$G$中的元素为基，即
    \begin{equation*}
        \mathbb{F}G = \ab\{\sum_{g \in G} a_g g \mid a_g \in \mathbb{F}\}.
    \end{equation*}
    我们定义其乘法为群乘法的双线性扩张，即对于任意$x = \sum_{g \in G} a_g g, y = \sum_{h \in G} b_h h \in \mathbb{F}G$，定义：
    \[ xy = \ab(\sum_{g \in G} a_g g) \ab(\sum_{h \in G} b_h h) = \sum_{g,h \in G} (a_g b_h) (gh) \]
    下面验证其为$\mathbb{F}-\!\!$代数：
    \begin{enumerate}
        \item \textbf{结合环}：由于群$G$的乘法满足结合律，且上述乘法定义继承了群$G$的结合律，因此$\mathbb{F}G$是结合环.
        \item \textbf{标量相容性}：对于任意$\lambda \in \mathbb{F}$，有\begin{equation*}\lambda(xy) = \lambda \sum_{g,h \in G} (a_g b_h) (gh) = \sum_{g,h \in G} (\lambda a_g) b_h (gh) = (\lambda x) y.
        \end{equation*}
        同理可得 $\lambda(xy) = \sum_{g,h \in G} a_g (\lambda b_h) (gh) = x(\lambda y)$，说明标量乘法与环乘法可交换.
    \end{enumerate}
    因此，$\mathbb{F}G$完全满足域上代数的条件.\par
    进一步，群表示$\rho:G \to \mathrm{GL}(V)$和群代数表示$\mathbb{F}G \to \mathrm{End}(V)$本质上是相同的.\par
    对于群代数中的任意元素$\sum_{i}a_ig_i$，可以自然地将基上的表示延拓到整个群代数，
    \begin{equation*}
        \left(\sum_{i}a_ig_i\right)v = \sum_{i}a_i\rho(g_i)v
    \end{equation*}
    反之，如果已知群代数的表示，那么只需要提取其基上的表示即可得到群表示.\par
    对于代数，我们也关心其中心，群代数的中心$Z(\mathbb{F}G)$的基为$\ab\{e_i = \sum_{g \in C_i} g\}$，其中$C_i$为$G$的共轭类.\index{z@中心!q@群代数}\index{q@群代数!z@中心}\par
    \begin{proof}
        设$x = \sum_{g \in G} a_g g \in Z(\mathbb{F}G)$，则$\forall\,h \in G$，有$hx = xh$，即$hxh^{-1} = x$.\par
    左式展开得$hxh^{-1} = \sum_{g \in G} a_g (hgh^{-1})$，令$g' = hgh^{-1}$，则$g = h^{-1}g'h$，于是$hxh^{-1} = \sum_{g' \in G} a_{h^{-1}g'h} g'$.\par
    将其与$x = \sum_{g' \in G} a_{g'} g'$比较系数可知，对任意$h, g' \in G$，有$a_{h^{-1}g'h} = a_{g'}$.\par
    这说明如果两个元素相互共轭，它们在中心元素$x$中的系数必定相同.\par
    因此，$x$可以提取出公共系数写为共轭类内部元素和的线性组合，即$x = \sum_i \lambda_i e_i$.\par
    另一方面，易验证对任意共轭类和$e_i$，都有$he_i h^{-1} = e_i$，即$e_i \in Z(\mathbb{F}G)$.\par
    又因为不同的共轭类对应集合互不相交，所以这些类和$\{e_i\}$是线性无关的，从而严格构成了$Z(\mathbb{F}G)$的基.
    \end{proof}
\end{instance}在了解代数的基本概念后，可以研究代数之间的结构.
若$A_1, A_2, \dots, A_n$为代数，可以通过它们之间定义运算.
\begin{align*}
    (a_1, \dots, a_n) + (b_1, \dots, b_n) &= (a_1 + b_1, \dots, a_n + b_n) \\
    \lambda(a_1, \dots, a_n) &= (\lambda a_1, \dots, \lambda a_n) \\
    (a_1, \dots, a_n)(b_1, \dots, b_n) &= (a_1b_1, \dots, a_nb_n)
\end{align*}
从而可对应定义代数直和直积.

为进一步研究代数的直和，这里引入幂等元的概念.

\begin{definition}[幂等元，中心幂等元，正交幂等元]\index{m@幂等元}\index{z@中心幂等元}\index{z@正交幂等元}
    设$A$是代数，若$e \in A$，且$e^2 = e$，则称$e$为\colorbf{幂等元}. \par
    特别地，若$e \in Z(A)$，称为\colorbf{中心幂等元}.\par
    若$e_1, e_2$均为幂等元，且$e_1 e_2 = e_2 e_1 = 0$，则称$e_1, e_2$为\colorbf{正交幂等元}.
\end{definition}

借助幂等元的概念，我们可以判断一个代数$A$能否分解为更小的代数结构.

\begin{proposition}[]
    设$A$是$F-\!\!$代数，则下述三者等价：
    \begin{enumerate}[(i)]
        \item $A$可分解成若干代数的直和，即$A = A_1 \times A_2 \times \dots \times A_n$
        \item $A$是其非平凡理想的直和，即$A = I_1 \oplus \dots \oplus I_n$，其中$I_i$为$A$的理想
        \item $A$的单位元可分解为两两正交的中心幂等元之和，即$1 = \sum\limits_i e_i$，$e_i^2 = e_i \in Z(A)$，$e_i e_j = 0\ (i \neq j)$
    \end{enumerate}
\end{proposition}

\begin{proof}
    (ii) $\Rightarrow$ (iii)：若$F-\!\!$代数$A$是其理想$I_1, \dots, I_n$的直和（作为向量空间的直和），则$A$的单位元$1$分解成两两正交的中心幂等元之和：$1 = e_1 + \dots + e_n$，其中$e_i \in I_i,\ i = 1, \dots, n$. 此时，$e_i$恰是$I_i$作为$F-\!\!$代数的单位元. \\
    (iii) $\Rightarrow$ (i)：设$A$的单位元$1$分解成两两正交的中心幂等元之和：$1 = e_1 + \dots + e_n$. 令$A_i := e_i A e_i$. 则$A_i$是$A$的$F-\!\!$子代数，且对任一$a \in A$有
    \begin{align*}
        a &= e_1 a + \dots + e_n a = e_1 (ae_1 + \dots + a e_n) + \dots + e_n(ae_1 + \dots + ae_n) \\
          &= \sum_{1 \le i, j \le n} e_i a e_j = e_1 a e_1 + \dots + e_n a e_n.
    \end{align*}
    于是得到$F-\!\!$代数同构$A \cong A_1 \times \dots \times A_n$. \\
    (i) $\Rightarrow$ (ii)：显然，将每个直积因子视为$A$的理想即得.
\end{proof}

\begin{proposition}[]
    设$A = A_1 \times A_2 \times \dots \times A_n$，$I$是$A$的任一理想，则$I = I_1 \oplus \dots \oplus I_n$，其中$I_i = I \cap A_i,\ 1 \le i \le n$. \\
    特别地，若$A_i$均为单代数，$I$是$A$的理想，则$I$必为若干$A_i$的直积.
\end{proposition}

\begin{proof}
    设$e_i$是$A_i$的单位元. 对于任一$x \in I$，有$x = x_1 + \dots + x_n$，其中$x_j \in A_j,\ 1 \le j \le n$. 
    因$e_i x_j = 0,\ j \neq i$，故$x_i = e_i x_i = e_i x \in I$. 
    这表明$I = I_1 \oplus \dots \oplus I_n$.
\end{proof}

\begin{definition}[本原幂等元，中心本原幂等元]\index{b@本原幂等元}\index{z@中心本原幂等元}
    幂等元$e$称为\colorbf{本原的}，如果$e$不能表成$A$的两个正交的幂等元的和.\par 
    幂等元$e$称为\colorbf{中心本原的}，如果$e \in Z(A)$且$e$不能表示成两个正交的中心幂等元的和.
\end{definition}

注意：中心本原的幂等元未必是本原的幂等元. 例如，若代数$A$是连通的（即$A$不能分解为两个非零代数的直积），则$A$的单位元是中心本原的幂等元，但却未必是本原的，即$1$仍可能分解成两个正交的幂等元之和.\par
若$e_1, \dots, e_n$是$A$的一组两两正交的幂等元，则$e_1, \dots, e_n$是$F-\!\!$线性无关的，因此$n \le \dim_F A$.

\begin{proposition}[]
    \begin{enumerate}[(i)]
        \item 设$e$是$A$的中心幂等元. 则$e$是中心本原的当且仅当$F-\!\!$代数$Ae$不能分解成$A$的非零理想的直和.
        \item 设$e_1, e_2$是$A$的两个中心本原幂等元. 则或$e_1 = e_2$，或$e_1$与$e_2$正交.
        \item 设$A$的单位元有分解$1 = e_1 + \dots + e_n$，且$e_1, \dots, e_n$是两两正交的中心本原幂等元（这相当于说$A$分解成连通的代数的直积）. 则$e_1, \dots, e_n$是$A$的全部中心本原幂等元.
    \end{enumerate}
\end{proposition}

\begin{proof}
    \begin{enumerate}[(i)]
        \item 类似于前面命题的推导.
        \item 我们有$e_1 = e_1 e_2 + e_1(1 - e_2)$. 因$e_1, e_2$均为中心幂等元，故可直接验证$e_1 e_2$与$e_1(1 - e_2)$均为中心幂等元并且互相正交. 又因$e_1$是中心本原的，故$e_1 e_2 = 0$，或$e_1(1 - e_2) = 0$，即$e_1 e_2 = 0$或$e_1 = e_1 e_2$.
        同理，我们有$e_1 e_2 = 0$或$e_2 = e_1 e_2$. 从而$e_1$与$e_2$正交或$e_1 = e_1 e_2 = e_2$.
        \item 假设$f$是$A$的中心本原幂等元，且$f \neq e_i,\ 1 \le i \le n$. 则由(ii)得到矛盾：$f = f e_1 + \dots + f e_n = 0$. \qedhere
    \end{enumerate}
\end{proof}

\section{Peirce 分解}

以上讨论了幂等元与代数直和分解的关系，下面将单位元的正交幂等元分解引向 Peirce 分解.

若 $A$ 的单位元 $1$ 可分解成两两正交的幂等元之和：$1=e_1+\cdots+e_n$，则左正则 $A-\!\!$模有直和分解 $A=Ae_1\oplus\cdots\oplus Ae_n$；（反之亦然；）同时，右正则 $A-\!\!$模也有直和分解 $A=e_1A\oplus\cdots\oplus e_nA$（反之亦然）. 于是得到 $F-\!\!$向量空间 $A$ 的直和分解.
\begin{definition}[Pierce分解]\index{p@Peirce 分解}
    设$A$是$F-\!\!$代数，$1=e_1+\cdots+e_n$是$A$的单位元的两两正交的幂等元分解. 则
    \begin{equation*}
        A = \bigoplus_{i,j=1}^n e_i A e_j
    \end{equation*}
    并称其为$A$的\colorbf{双边Peirce分解}.
\end{definition}

注意，一般来说，$e_iAe_j$ 既非 $A$ 的左理想，也非 $A$ 的右理想；但是此分解却允许我们将 $A$ 的元素表达成矩阵形式：
\[
\forall\, a\in A, \quad a=\begin{pmatrix} a_{11} & \cdots & a_{1n} \\ \vdots & \ddots & \vdots \\ a_{n1} & \cdots & a_{nn} \end{pmatrix},
\]
其中 $a_{ij}=e_iae_j\in e_iAe_j$；而且 $A$ 中元素的运算可归结为相应矩阵的运算，例如
\[
ab=\sum_{i,k}\sum_{l,j} a_{ik}b_{lj} = \sum_{i,j}\left(\sum_k a_{ik}b_{kj}\right),
\]
这是因为当 $k\ne l$ 时，$a_{ik}b_{lj}=e_ia(e_k\cdot e_l)be_j=0$. 由此可以将 $A$ 写成矩阵形式
\[
A=\begin{pmatrix} A_{11} & \cdots & A_{1n} \\ \vdots & \ddots & \vdots \\ A_{n1} & \cdots & A_{nn} \end{pmatrix},
\]
其中 $A_{ij}=e_iAe_j$.

代数上的 Peirce 分解自然引出模的直和分解与自同态代数幂等元分解之间的对应.

设模 $M$ 有直和分解 $M=M_1\oplus\cdots\oplus M_n$. 令 $e_j: M\longrightarrow M_j$ 是投影映射. 则 $e_j$ 可自然地看成 $\End_A(M)$ 中元，于是 $\End_A(M)$ 的单位元分解为两两正交的幂等元之和：$1=e_1+\cdots+e_n$；反之，令 $M_i=e_i(M)$. 则 $M_i$ 是左 $A$-模，易知 $M=M_1\oplus\cdots\oplus M_n$. 由此即得

\begin{proposition}[]
    模 $M$ 的直和分解与代数 $\End_A(M)$ 的单位元的正交幂等元分解是一一对应的.
\end{proposition}

称模 $M$ \colorbf{可分解}\index{k@可分解模}，如果存在 $M$ 的两个非零子模 $M_1$ 和 $M_2$ 使得 $M=M_1\oplus M_2$；否则，称 $M$ \colorbf{不可分解}\index{b@不可分解模}. 因此，若 $M$ 不可分解且 $M=M_1\oplus M_2$，则必推出 $M_1=0$ 或 $M_2=0$.

\begin{corollary}[推论 1]\index{b@不可分解模!d@等价刻画}
    非零模 $M$ 不可分解当且仅当 $\End_A(M)$ 仅有两个平凡的幂等元：$0$ 和 $1$.
\end{corollary}

\begin{proof}
    若 $e$ 是 $\End_A(M)$ 非平凡的幂等元，则 $f=1-e$ 也是非平凡的幂等元且 $ef=fe=0$，$1=e+f$. 从而由上述命题 $M$ 是可分解的.
\end{proof}

根据 2.9 中的讨论，代数 $\End_A(M)$ 的单位元的正交幂等元分解与左正则 $\End_A(M)$-模的直和分解是一一对应的. 由此即得

\begin{corollary}[推论 2]\index{z@直和分解!d@与左正则模的对应}
    左 $A$-模 $M$ 的直和分解与代数 $\End_A(M)$ 的左正则模的直和分解是一一对应的.
\end{corollary}

% 过渡：利用模直和分解对应的幂等元结构，自同态代数中的元素可获得矩阵表示.

2.11 设左 $A$-模 $M$ 有直和分解 $M=M_1\oplus\cdots\oplus M_n$. 则 $\End_A(M)$ 的单位元 1 有分解 $1=e_1+\cdots+e_n$，其中 $e_i: M\longrightarrow M_i$ 是投射. 对于 $f\in\End_A(M)$，令 $f_{ji}: M_i\longrightarrow M_j$ 是左 $A$-模同态，其中 $f_{ji}(x)=e_jf(x)$，$\forall\, x\in M_i$.
设 $m=m_1+\cdots+m_n\in M$，$m_i\in M_i$，$i=1,\cdots,n$. 则 $f(m)=f(m_1)+\cdots+f(m_n)$. 设 $f(m_i)=m_{i1}+\cdots+m_{in}\in M$，其中 $m_{ij}\in M_j$，$j=1,\cdots,n$. 于是
\[
f(m)=\sum_{1\le i,j\le n} m_{ij} = \sum_{1\le i,j\le n} f_{ji}(m_i).
\]
因此，若将 $m$ 写成列向量 $\begin{pmatrix} m_1 \\ \vdots \\ m_n \end{pmatrix}$，将 $f$ 写成矩阵形式 $f=(f_{ij})_{n\times n}$，则
\[
f(m)=\begin{pmatrix} f_{11} & \cdots & f_{1n} \\ \vdots & \ddots & \vdots \\ f_{n1} & \cdots & f_{nn} \end{pmatrix}\begin{pmatrix} m_1 \\ \vdots \\ m_n \end{pmatrix}.
\]
这种记号的好处是 $\End_A(M)$ 中元素 $f$ 的加法、乘法和数乘与其相应的矩阵 $(f_{ij})_{n\times n}$ 的加法、乘法和数乘是完全一致的（请读者自行验证）. 于是，我们得到

\begin{lemma}[]
    有 $F$-代数同构
    \[
    \End_A(M_1\oplus\cdots\oplus M_n) \cong \left\{ \begin{pmatrix} f_{11} & \cdots & f_{1n} \\ \vdots & \ddots & \vdots \\ f_{n1} & \cdots & f_{nn} \end{pmatrix} \Bigg| f_{ij}\in\Hom_A(M_j, M_i),\; 1\le i, j\le n \right\},
    \]
    其中右边集合的代数运算按矩阵运算法则进行.
\end{lemma}

特别地，若 $\Hom_A(M_i, M_j)=0$，$1\le i\ne j\le n$，则
\[
\End_A(M_1\oplus\cdots\oplus M_n) \cong \End_A(M_1)\times\cdots\times\End_A(M_n).
\]

在上述引理中取 $M_1\cong\cdots\cong M_n\cong L$，即 $M\cong nL$，则得到

\begin{corollary}[]
    $\End_A(nL)\cong M_n(B)$，其中 $B=\End_A(L)$，$M_n(B)$ 表示代数 $B$ 上的全矩阵代数.
\end{corollary}

\section{模同态}

\begin{definition}[模同态]
    设$M, M'$均为左$A-\!\!$模，$f: M \to M'$是群同态，若\begin{equation}
    \forall\, a \in A, m \in M, f(am) = a f(m),
    \end{equation}
    则称$f$是\colorbf{模同态}.
\end{definition}

定义模同态后就可以照旧研究其核和像.除此之外，还关注以下两个空间.

\begin{definition}[余核，余像]{def:余核}\index{y@余核}\index{y@余像}
    设$M, M'$是左$A-\!\!$模，$f: M \to M'$是模同态，则定义其余核与余像分别为：
    \begin{align*}
        \coker f &\triangleq M' / \im f \\
        \coim f &\triangleq M / \ker f
    \end{align*}
\end{definition}

关于单同态和满同态，有若干有用的等价条件.

\begin{proposition}[单射的等价条件]\index{d@单模同态}
    设$f: M \to M'$是模同态，则以下等价：
    \begin{enumerate}[(i)]
        \item $f$是单射
        \item $\ker f = 0$
        \item $\forall\, g, h: N \to M$，$fg = fh \Rightarrow g = h$
        \item $\forall\, g: N \to M$，$fg = 0 \Rightarrow g = 0$
    \end{enumerate}
\end{proposition}

\begin{proposition}[满射的等价条件]\index{d@满模同态}
    设$f: M \to M'$是模同态，则以下等价：
    \begin{enumerate}[(i)]
        \item $f$是满射
        \item $\mathrm{Im}\, f = M'$
        \item $\forall\, g, h: M' \to N$，$gf = hf \Rightarrow g = h$
        \item $\forall\, g: M' \to N$，$gf = 0 \Rightarrow g = 0$
    \end{enumerate}
    
\end{proposition}

类似于群同态的情况，模同态也可作分解.

\begin{theorem}[][thm:模同态分解]
    设$f: M \to M'$是模同态，$g: M \to N$为满射，且$\ker g \subseteq \ker f$. 则下述交换图成立：
    \begin{center}
    \begin{tikzcd}[labels={font=\normalsize}]
        M \arrow[rr, "f"] \arrow[dd, "g"'] & & M' \\
        \\
        N \arrow[rruu, "\exists !\, h"', dashed]
    \end{tikzcd}
    \end{center}
    即$\exists\, !\, h: N \to M',\st f = hg$.\par
    特别地，$h$是单射 $\Leftrightarrow \ker g = \ker f$.
    $h$是满射 $\Leftrightarrow f$是满射.
\end{theorem}

上述定理是在证明模同态时关键的工具，下面是一些重要的结论

\begin{theorem}[]
    设$f: M \to M'$是群同态，则
    \begin{enumerate}[(i)]
        \item $M / \ker f \cong M'$ 当且仅当$f$是满射
        \item 若$K \subseteq N \subseteq M$，$M / N \cong (M / K) / (N / K)$.
        \item $(N + K) / K \cong N / (N \cap K)$.
    \end{enumerate}
\end{theorem}

\begin{proof}
    \begin{enumerate}[(i)]
        \item 首先有典范满同态$g: M \twoheadrightarrow M / \ker f$.\par
        由\cref{thm:模同态分解}，由于$\ker g = \ker f$，$\exists\, h: M / \ker f \to M',\st f = hg$.\par
        
        若$f$满，则$h$满.
        又$\ker g = \ker f$，则$h$单.
        故$h$为同构，$M / \ker f \cong M'$.\par
        反之，若$M / \ker f \cong M'$，则由$h, g$满得$f$满.
        \item 有典型满同态$f: M/K \twoheadrightarrow (M/K)/(N/K)$.\par
        由于$K \subseteq N$，可定义满射
        \begin{align*}
            g: M/K &\to M/N, \\[-7pt]
            x+K &\mapsto x+N.
        \end{align*}
        若$g(x+K) = 0$，则$x \in N$且$x \notin K$，故$\ker g = N/K$.\par
        由于$\ker f = N/K = \ker g$，故$\exists\, h: M/N \to (M/K) / (N/K)$.\par
        由$\ker f = \ker g$，$h$单.由$f, g$满，知$h$满.因此$M/N \cong (M/K) / (N/K)$.
        \item 同理，考虑$N \to (N+K)/K$即可. \qedhere
    \end{enumerate}
    \begin{figure}[H]
        \centering
        \renewcommand{\thesubfigure}{\roman{subfigure}}
        \hspace{-0.1\textwidth}
        \begin{subfigure}[b]{0.2\linewidth}
            \centering
            \begin{tikzcd}[labels={font=\normalsize}]
                M \arrow[r, "f"] \arrow[dd, "g"',two heads] & M' \\
                \\
                M / \ker f \arrow[ruu, "h"', dashed]
            \end{tikzcd}
        \end{subfigure}
        \hspace{0.05\linewidth}
        \begin{subfigure}[b]{0.28\linewidth}
            \centering
            \begin{tikzcd}[labels={font=\normalsize}]
                M/K \arrow[r, "f",two heads] \arrow[dd, "g"'] & (M/K)/(N/K) \\
                \\
                M/N \arrow[ruu, "h"', dashed]
            \end{tikzcd}
        \end{subfigure}
        \hspace{0.05\linewidth}
        \begin{subfigure}[b]{0.28\linewidth}
            \centering
            \begin{tikzcd}[labels={font=\normalsize}]
                N \arrow[r, "f",two heads] \arrow[dd, "g"'] & N/(N\cap K) \\
                \\
                (N+K)/K \arrow[ruu, "h"', dashed]
            \end{tikzcd}
        \end{subfigure}
        \caption{证明示意.}
    \end{figure}
\end{proof}

下面我们考虑模同态序列.

\begin{definition}[正合，正合列][]\index{z@正合列}
    设$M, M', M''$为模，且有模同态列
    \begin{equation*}
        \exactseq{M'}[f]{M}[g]{M''}
    \end{equation*}
    若$\mathrm{Im}\, f = \ker g$，则称其在$M$ \colorbf{正合}.\par
    若序列
    \begin{equation*}
        \exactseq{\cdots}[]{M_{i-1}}[]{M_i}[]{M_{i+1}}[]{\cdots}
    \end{equation*}
    在每一个$M_i$处正合，则称其为\colorbf{正合列}.
\end{definition}

\begin{instance}
    \begin{enumerate}
        \item 若\exactseq{0}[0]{M'}[f]{M}正合，$\ker f = \im 0 = 0$，故$f$是单射.
        \item 若\exactseq{M}[g]{M''}[0]{0}正合，则$\mathrm{Im}\, g = \ker 0 = M''$，故$g$是满射.
    \end{enumerate}
\end{instance}

由此可以解释为何要定义余核与余像.
假设有以下缺失的正合列：
\begin{equation*}
    \exactseq{0}[0]{\ker f}[]{M}[f]{N}[g]{?}[]{0}
\end{equation*}
为使其为正合列，有$g$是满射，且$\ker g = \mathrm{Im}\, f$，从而问号处应为$N / \mathrm{Im}\, f$，这就是我们定义的余核.\par
特别地，若上述$f$单，则$\ker f = 0$. 此时两个序列合二为一，由正合列变为
\begin{equation*}
    \exactseq{0}[]{M}[f]{N}[]{\coker f}[]{0}
\end{equation*}
若$f$满，则$N / \mathrm{Im}\, f = 0$，故右边两个序列合二为一，此时正合列为
\begin{equation*}
    \exactseq{0}[]{\ker f}[]{M}[f]{N}[]{0}
\end{equation*}

这两个正合列都只包含三个非零的模，除了例3.3外的情况外，它们是最短的正合列，因此被称为\colorbf{短正合列}.

\begin{definition}[短正合列][def:短正合列]\index{d@短正合列}
    \colorbf{短正合列}为以下形式的正合列：
    \begin{equation*}
        \exactseq{0}[]{M_1}[f]{M_2}[g]{M_3}[]{0}
    \end{equation*}
\end{definition}
事实上，可以将短正合列理解为模$M$与其子模$K$及商模$M/K$组成的正合列：
\[\exactseq{0}[]{K}[]{M}[f]{M/K}[]{0}\]

研究短正合列的另一个原因是所有的正合列都可以由短正合列拼接而成.这一点将在最后范畴的学习中进一步说明.

\section{半单模与半单代数}

了解了模同态和短正合列后，我们希望通过最简单的模块来研究复杂的模.为此，我们引入半单模的概念，它类似于向量空间可以分解为一维子空间的直和.

\begin{definition}[半单模]
    设$M$是$A-\!\!$模，则称$M$为\colorbf{半单模}，若其满足下述任一等价条件：
    \begin{enumerate}[(i)]
        \item $M = \bigoplus M_i$，其中$M_i$是单模
        \item $\forall\, N \subseteq M, \exists\, W, \text{s.t. } M = N \oplus W$
        \item $M$为其单子模之和
    \end{enumerate}
\end{definition}

由此可获得Schur引理推广到模上的情形.

\begin{lemma}[Schur引理]
    设$M, N$均为单$A-\!\!$模，则
    \begin{enumerate}[(i)]
        \item 若$M \not\cong N$，则$\mathrm{Hom}_A(M, N) = 0$
        \item 若$M \cong N$，$\mathrm{End}_A(M) = \mathrm{Hom}_A(M, M) = \mathrm{Hom}_A(M, N)$是可除代数.
        特别地，若$F$是代数闭域，则$\mathrm{Hom}_A(M, N) \cong F I_M \cong F$.
    \end{enumerate}
\end{lemma}

\begin{definition}[半单代数]
    设$A$为代数，若任意$A-\!\!$模$M$都为半单模，则称$A$为\colorbf{半单代数}.
\end{definition}

对于上述定义，一个经典的例子是有限群群代数的半单性：

\begin{instance}
    $F S_n$是半单代数 $\Leftrightarrow \mathrm{char}\, F > n$ (或 $\mathrm{char}\, F = 0$).
\end{instance}

\section{Jordan-Hölder 定理}

本节要证明有限维代数的有限维表示中的两条基本定理.

\textbf{3.1} 设 $A$ 是有限维 $F-\!\!$代数， $M$ 是有限维 $A-\!\!$模. $M$ 的子模链
\[ M = M_0 \supset M_1 \supset \cdots \supset M_{l-1} \supset M_l = 0 \]
称为长度为 $l$ 的\colorbf{合成列}\index{h@合成列}，如果所有商模 $M_i/M_{i+1}$ 均为单 $A-\!\!$模， $0 \le i \le l-1$. 此时， $M_i/M_{i+1}$ 称为该合成列的\colorbf{合成因子}\index{h@合成因子}.

任一（有限维）模 $M$ 总有合成列. 事实上，若 $M$ 单，则 $M \supset 0$ 是合成列；若 $M$ 非单，则 $M$ 有极大子模 $M_1$（即不存在子模 $X$ 使得 $M \supsetneq X \supsetneq M_1$. 因 $\dim_F M < +\infty$，故 $M_1$ 总存在）. 若 $M_1$ 单，则 $M \supset M_1 \supset 0$ 是合成列；若 $M_1$ 非单，则重复上述过程. 因 $M$ 有限维，故最后总得到 $M$ 的一个合成列.

若 $M/N \cong L$，则称 $M$ 是 $N$ 借助 $L$ 的\colorbf{扩张}\index{k@扩张!m@模的扩张}. 模 $M$ 有合成列这一事实说明：从 $M$ 的单子模出发，通过借助单模的反复扩张可得到 $M$.

模 $M$ 的合成列一般来说不唯一. 但下述定理表明对于任一单模 $S$，$M$ 的合成列中与 $S$ 同构的合成因子的个数不依赖于合成列的选取. 从而将 $M$ 的一个合成列的合成因子称为 $M$ 的合成因子.

\begin{theorem}[Jordan-Hölder 定理]\index{J@Jordan-Hölder 定理}
设 $M$ 是任一 $A-\!\!$模，$M = M_0 \supset M_1 \supset \cdots \supset M_s = 0$ 和 $M = N_0 \supset N_1 \supset \cdots \supset N_t = 0$ 是 $M$ 的两个合成列. 则 $s = t$，并且存在 $\{1,\cdots,s\}$ 的一个置换 $\pi$ 使得 $N_i/N_{i+1} \cong M_{\pi(i)}/M_{\pi(i)+1}$， $0 \le i \le s-1$.
\end{theorem}

\begin{proof}
对 $s$ 用归纳法. 结论对单 $A-\!\!$模成立. 假设定理对于具有长为 $s-1$ 的合成列的模成立. 设 $M$ 具有长为 $s$ 的合成列，形如题设.

若 $M_1 = N_1$. 则 $M_0/M_1 = N_0/N_1$，结论由归纳假设即得.

设 $M_1 \neq N_1$. 则 $M_1 + N_1 \neq M_1$. 因 $M_1$ 是 $M$ 的极大子模，故 $M_1 + N_1 = M$.
从而有 $A-\!\!$模同构（引理 2.2(ii)）
\[ N_1/(M_1 \cap N_1) \cong M/M_1; \quad M_1/(M_1 \cap N_1) \cong M/N_1. \]

取定 $M_1 \cap N_1$ 的一个合成列 $M_1 \cap N_1 = L_2 \supset L_3 \supset \cdots \supset L_k = 0$. 则
$M_1 \supset M_2 \supset \cdots \supset M_s = 0$ 与 $M_1 \supset L_2 \supset \cdots \supset L_k = 0$ 均为 $M_1$ 的合成列，由归纳假设知它们等价，即 $s = k$，且这两个合成列的合成因子之间存在一一对应，使得相对应的合成因子同构. 又因为 $N_1 \supset N_2 \supset \cdots \supset N_t = 0$ 与 $N_1 \supset L_2 \supset \cdots \supset L_s = 0$ 均为 $N_1$ 的合成列，由归纳假设知它们等价，即 $s = t$，且这两个合成列的合成因子之间存在一一对应.

现在考虑 $M$ 的合成列 $M \supset M_1 \supset L_2 \supset \cdots \supset L_s = 0$ 与 $M \supset N_1 \supset L_2 \supset \cdots \supset L_s = 0$. 两者的合成因子从第三个开始全部相同，而
\[ M/M_1 \cong N_1/L_2, \quad M_1/L_2 \cong M/N_1. \]
因此这两个合成列等价.

由此看出，合成列 $M \supset M_1 \supset \cdots \supset M_s = 0$ 等价于 $M \supset M_1 \supset L_2 \supset \cdots \supset L_s = 0$，又等价于 $M \supset N_1 \supset L_2 \supset \cdots \supset L_s = 0$，从而等价于 $M \supset N_1 \supset \cdots \supset N_s = 0$.
\end{proof}

\textbf{注}. 将单模 $S$ 出现在 $M$ 的合成因子中的次数记为 $d_S$，因此得到 $M$ 的维数向量 $(d_S)_S$. 这是模的一个重要不变量. 维数向量相同的两个不可分解模何时同构是代数表示论中有趣的问题.

\begin{theorem}[]\index{d@定理!d@单模的同构类有限性}
设 $A$ 是有限维代数. 则任一单 $A-\!\!$模 $S$ 均同构于正则模 $A$ 的某一合成因子. 特别地，仅有有限个单 $A-\!\!$模的同构类.
\end{theorem}

\begin{proof}
取左正则 $A-\!\!$模 ${}_A A$ 的一个合成列 $A = M_0 \supset M_1 \supset M_2 \supset \cdots \supset M_t = 0$.
则有 $A-\!\!$模的短正合列
\begin{align*}
0 &\longrightarrow M_1 \longrightarrow A \longrightarrow A/M_1 \longrightarrow 0, \\
0 &\longrightarrow M_2 \longrightarrow M_1 \longrightarrow M_1/M_2 \longrightarrow 0, \\
&\cdots, \\
0 &\longrightarrow M_{t-1} \longrightarrow M_{t-2} \longrightarrow M_{t-2}/M_{t-1} \longrightarrow 0.
\end{align*}
利用左正合函子 $\Hom_A(-, S)$（参见习题 2.10），我们得到 $F-\!\!$模正合列
\begin{align*}
0 &\longrightarrow \Hom_A(A/M_1, S) \longrightarrow \Hom_A(A, S) \longrightarrow \Hom_A(M_1, S), \\
0 &\longrightarrow \Hom_A(M_1/M_2, S) \longrightarrow \Hom_A(M_1, S) \longrightarrow \Hom_A(M_2, S), \\
&\cdots, \\
0 &\longrightarrow \Hom_A(M_{t-2}/M_{t-1}, S) \longrightarrow \Hom_A(M_{t-2}, S) \longrightarrow \Hom_A(M_{t-1}, S).
\end{align*}
因为 $\Hom_A(A, S) \cong S \neq 0$，由上述正合列可推出存在 $0 \le i \le t-1$，使得 $\Hom_A(M_i/M_{i+1}, S) \neq 0$. 于是，由 Schur 引理知 $S \cong M_i/M_{i+1}$.
\end{proof}

\section{自由模}
下面我们将介绍一系列特殊的模.首先是自由模，它可以视作是向量空间在模上的推广.\par

\begin{definition}[自由模]\index{z@自由模}
    设$M$是$R-\!\!$模，若存在$\{a_i\}_{i\in I} \subset M$，且满足：
    \begin{enumerate}[(i)]
        \item $\forall\, x \in M$, $x = \sum_{i \in I} r_i a_i, r_i \in R$
        \item 若$\sum_{i \in I} r_i a_i = 0 \Rightarrow \forall\, r_i = 0$.
    \end{enumerate}
    则称$M$是\colorbf{自由模}，$\{a_i\}_{i\in I}$是$M$的一组\colorbf{基}.
\end{definition}

以下定理可以帮助我们理解自由模的结构.
\begin{proposition}[]
    $M$是自由模当且仅当$ M \cong \bigoplus M_i, M_i \cong R$.
\end{proposition}
\begin{proof}
    \fbox{$\Rightarrow$} 取$M_i = R a_i$，则$M_i \subset M$且$M_i \cong R$.
    由于$\{a_i\}_{i\in I}$为$M$的基，$M = \bigoplus M_i$.

    \fbox{$\Leftarrow$} 设$\varphi: \bigoplus_{i \in I} M_i \to M$为同构映射，令$f_i: R \to M_i$为每个分量的同构.\par
    设$1_R$为$R$的单位元，令$a_i = \varphi(f_i(1_R))$. \par
    对于任意$x \in M$，由于$\varphi$是同构，存在唯一标量族$r_i \in R$使得$\varphi^{-1}(x) = (f_i(r_i))_{i\in I}$.\par
    因此\begin{equation*}
        x = \sum_{i \in I} r_i \varphi(f_i(1_R)) = \sum_{i \in I} r_i a_i.
    \end{equation*}

    且若$\sum_{i \in I} r_i a_i = 0$，将$\varphi^{-1}$作用于两边，利用直和的性质得$f_i(r_i) = 0$，进而由$f_i$同构得$r_i = 0$. \par
    因此$\{a_i\}_{i \in I}$为$M$的基，$M$是自由模.
\end{proof}

若$M$是自由模，$X$为其基，则若有$X \to N$的模映射$f$，有如下交换图成立：
\begin{center}
\begin{tikzcd}
    X \arrow[r, "f"] \arrow[d, hook, "\iota"'] & N \\
    M \arrow[ru, "\exists\, !\, \bar{f}"', dashed] &
\end{tikzcd}
\end{center}
这里的$\bar{f}$可以显示地构造出来\begin{equation*}
    \bar{f}(x) = \bar{f}\ab(\sum r_i a_i) = \sum r_i f(a_i).
\end{equation*}
和线性空间类似，这告诉我们想确定自由模的模同态，只需确定其在基上的值即可.\par
下一个讨论的内容是自由群的「自由」体现在何处.在群论中，任何群都可以表为自由群的商群.这一点对自由模同样成立.

\begin{proposition}[]
    对任意$R-\!\!$模$M$，都存在自由$R-\!\!$模$V$，使得$M \cong V/N$.
\end{proposition}

\begin{proof}
    考虑$M$的生成元集$X$，即$\forall\, x \in M$, $x = \sum_{x_i \in X} r_i x_i$. 定义$V \triangleq \bigoplus_{x_i \in X} R x_i$.\par
    记$\iota: X \to M$为嵌入. 定义\begin{align*}
        \bar{f}: V &\to M \\
        x &\mapsto \sum a_i \iota(x_i)
    \end{align*}
    显然$\bar{f}$满，故$M \cong V / \ker \bar{f}$.
\end{proof}

\begin{proposition}[]
    设$V$是自由模，$g: V \to N$且存在满同态$f: M \to N$，则下述交换图成立：
    \begin{center}
    \begin{tikzcd}
        V \arrow[r, "g"] \arrow[d, "\exists\, \bar{h}"', dashed] & N \\
        M \arrow[ru, "f"'] &
    \end{tikzcd}
    \end{center}
    即$f \bar{h} = g$.
\end{proposition}

\begin{proof}
    只需定义$\bar{h}(x_i)$，使得$f(\bar{h}(x_i)) = g(x_i)$.\par
    由于$f$是满射，$\exists\, m_i \in M, \st f(m_i) = g(x_i)$.\par
    令$\bar{h}(x_i) = m_i$即可.
\end{proof}

自由模作为线性空间的推广，其部分性质没有线性空间好.以下面两个定理为例，线性空间的基中元素个数必相等，但这一性质对更广泛的自由模并不成立.

\begin{theorem}[]\index{z@自由模!j@秩}
    若$R$是交换环，$M$为自由$R-\!\!$模，则$M$的任一基具有相同个数，称其为$M$的秩，记作$\operatorname{rank} M$.
\end{theorem}

另一个例子线性空间的子空间仍是线性空间，并且维数一定小于等于原空间.但自由模的子模不一定是自由模，并且即使是自由模的子模也是自由模时，其秩也不一定小于等于原空间的秩.

\begin{theorem}[]\index{z@自由模!z@子模的自由性}
    设$V$是自由$R-\!\!$模，则任意$V$的子模$M$是自由$R-\!\!$模且$\operatorname{rank} M \leq \operatorname{rank} V \Leftrightarrow R$是主理想整环.
\end{theorem}
\begin{proof}
    $\Rightarrow$ 若取$V=R$，则$V$是秩为$1$的自由$R-\!\!$模. $R$的任何理想$I$都是$V$的子模. 
    根据假设，$I$是秩为$0$或$1$的自由模. 若$I=0$，则即为主理想$(0)$；
    若$I \neq 0$，则$I$的秩为$1$，故$I \cong R$，即存在$a \in I$使得$I = Ra$，故$I$是主理想.
    且因为$I \cong R$，同构映射$R \to Ra, r \mapsto ra$是单射，这意味着如果$ra=0$必有$r=0$（其中$a \neq 0$），进而推出$R$没有零因子，所以$R$是整环. 因此$R$是主理想整环.

    $\Leftarrow$ 设$R$是主理想整环，设$V$是自由$R-\!\!$模，基为$\{e_i\}_{i \in I}$. 利用良序定理将指标集$I$良序化. 对任意$j \in I$，令$V_j$为由$\{e_i \mid i < j\}$生成的子模，$\bar{V}_j$为由$\{e_i \mid i \le j\}$生成的子模. 
    对$V$的任一子模$M$，考虑投影同态$\pi_j: M \cap \bar{V}_j \to R$，其定义是将元素映射到其在$e_j$上的坐标. $\pi_j$的像$\im \pi_j$是$R$的一个理想，由于$R$为主理想整环，该理想由某个$a_j \in R$生成.
    若$a_j = 0$，则令$B_j = \emptyset$；若$a_j \neq 0$，取$x_j \in M \cap \bar{V}_j$使得$\pi_j(x_j) = a_j$，此时令$B_j = \{x_j\}$.
    容易验证正交集合$B = \bigcup_{j \in I} B_j$是$M$的一个独立生成集（一组基）.并且，存在自然单射将$B$的元素对应到指标集$I$中那些$a_j \ne 0$的子集，从而$M$是自由模，且$\operatorname{rank} M = |B| \le |I| = \operatorname{rank} V$.
\end{proof}
\section{内射模和投射模}
在了解了模的基本同态和单模、半单模性质后，我们回过头来看模的直和.
类比于其它代数结构，若有$R-\!\!$模$M_1, M_2$，有其直和$M_1 \oplus M_2$，此时$\forall\, r \in R$，$r(m_1, m_2) \triangleq (rm_1, rm_2)$.
关于直和有典范内射
\begin{align*}
    \iota_1: M_1 &\to M_1 \oplus M_2 & \iota_2: M_2 &\to M_1 \oplus M_2 \\[-7pt]
    x_1 &\mapsto (x_1, 0) & x_2 &\mapsto (0, x_2)
\end{align*}
以及投影
\begin{align*}
    \pi_1: M_1 \oplus M_2 &\to M_1 & \pi_2: M_1 \oplus M_2 &\to M_2 \\[-7pt]
    (x_1, x_2) &\mapsto x_1 & (x_1, x_2) &\mapsto x_2
\end{align*}

利用上述四个映射可以构造出恒等映射：
\[ \pi_1 \iota_1 = \bm{1}_{M_1}, \quad \pi_2 \iota_2 = \bm{1}_{M_2}, \quad \pi_1 \iota_2 = \pi_2 \iota_1 = 0, \quad \iota_1 \pi_1 + \iota_2 \pi_2 = \bm{1}_{M_1 \oplus M_2}. \]
并且有以下$2$个短正合列：
\begin{align*}
    \exactseq{0}[]{M_1}[\iota_1]{M_1 \oplus M_2}[\pi_2]{M_2}[]{0} \\
    \exactseq{0}[]{M_2}[\iota_2]{M_1 \oplus M_2}[\pi_1]{M_1}[]{0}
\end{align*}

考察$f: N \to M$，$g: M \to N$，且$gf = \bm{1}_N$，则可知$f$是单射，$g$是满射，并且$M = \im f \oplus \ker g$.
后者的原因是$\forall\, x \in M$，$x = fg(x) + (x - fg(x))$，其中$fg(x) \in \im f$，$g(x - fg(x)) = g(x) - gfg(x) = g(x) - \bm{1}_N g(x) = 0$，故$x - fg(x) \in \ker g$.我们称$f$为分裂单的，$g$为分裂满的.
\begin{definition}[分裂单同态、分裂满同态][]\index{f@分裂单同态}\index{f@分裂满同态}
    设$M, N$是模，$f: N \to M$是模同态，若$\exists\, f': M \to N, \st f'f = \bm{1}_N$，则称$f$为\colorbf{分裂单同态}.\par
    设$M, N$是模，$g: M \to N$是模同态，若$\exists\, g': N \to M, \st gg' = \bm{1}_N$，则称$g$为\colorbf{分裂满同态}.
\end{definition}

利用这种同态的分裂性质，我们可以考察更特殊的短正合列，这类短正合列相当于给出了模的直和分解的另一种视角.

\begin{definition}[分裂正合列]\index{f@分裂正合列}
    若短正合列
    \[ \exactseq{0}[]{M_1}[f]{M_2}[g]{M_3}[]{0} \]
    满足$f$是分裂单的，$g$是分裂满的，则称该正合列为\colorbf{分裂正合列}.
\end{definition}

对于分裂正合列，可以证明$M_2 \cong M_1 \oplus M_3$.
实际上，正合列不要求两个分裂，它有以下若干等价条件.

\begin{proposition}[分裂正合列的等价条件]\index{f@分裂正合列}
    设$\exactseq{0}[]{M_1}[f]{M_2}[g]{M_3}[]{0}$是短正合列，则下述条件等价：
    \begin{enumerate}[(i)]
        \item 上述正合列为分裂正合列.
        \item $f$分裂单.
        \item $g$分裂满.
        \item $M_2 = \im f \oplus M' = \ker g \oplus M'$.
        \item $\forall\, h: M_1 \to N$，$ \exists\, \overline{h}: M_2 \to N$使交换图成立
        \begin{center}
        \begin{tikzcd}[cramped]
            0 \arrow[r] & M_1 \arrow[r, "f"] \arrow[d, "h"'] & M_2 \arrow[dl, "\overline{h}", dashed] \\
            & N
        \end{tikzcd}\end{center}
        即$\overline{h} f = h$.
        \item $\forall\, h: N \to M_3$，$ \exists\, \overline{h}: N \to M_2$使交换图成立
        \begin{center}\begin{tikzcd}[cramped]
            M_2 \arrow[r, "g"] & M_3 \arrow[r] & 0 \\
            N \arrow[ru, "h"'] \arrow[u, "\overline{h}", dashed]
        \end{tikzcd}\end{center}
        即$g \overline{h} = h$.
        \item 存在$\Phi:M_2\to M_1\otimes M_3$，使得交换图成立\par
        \begin{center}\begin{tikzcd}[cramped,column sep=scriptsize, ampersand replacement=\&]
            0 \arrow[r] \& M_1 \arrow[r, "f"] \arrow[d, "\bm{1}_{M_1}"'] \& M_2 \arrow[r, "g"] \arrow[d, "\cong"', "\Phi"] \& M_3 \arrow[r] \arrow[d, "\bm{1}_{M_3}"'] \& 0 \\
            0 \arrow[r] \& M_1 \arrow[r, "\iota_1"'] \& M_1 \oplus M_3 \arrow[r, "\pi_2"'] \& M_3 \arrow[r] \& 0
        \end{tikzcd}\end{center}
    \end{enumerate}
\end{proposition}

\begin{proof}
    证明思路为 $\text{(i)} \Leftrightarrow \text{(ii)} \Leftrightarrow \text{(vii)}$，$\text{(ii)} \Leftrightarrow \text{(iv)} \Leftrightarrow \text{(v)}$，$\text{(iii)} \Leftrightarrow \text{(vi)} \Leftrightarrow \text{(vi)}$.
    
    (i) $\Rightarrow$ (ii), (vii) $\Rightarrow$ (i) 显然.
    
    (ii) $\Rightarrow$ (iv). 由分裂单定义，$\exists\, f': M_2 \to M_1, \st f'f = \bm{1}_{M_1}$，则$M_2 = \im f \oplus \ker f'$.
    
    (iv) $\Rightarrow$ (v). 已知$M_2 = \im f \oplus M'$，定义
    \begin{align*}
    \overline{h}: \quad\qquad M_2 &\to N \\[-7pt]
    f(m_1) + m' &\mapsto h(m_1)
    \end{align*}
    则$\overline{h} f(m_1) = h(m_1)$，满足要求.
    
    (v) $\Rightarrow$ (ii). 取$h = \bm{1}_{M_1}$，则$\overline{h} f = \bm{1}_{M_1}$，故$f$分裂单.
    
    (iii) $\Leftrightarrow$ (iv) $\Leftrightarrow$ (vi) 同理可证.
    
    (ii) $\Rightarrow$ (vii). 由于$f$分裂单，$\exists\, h: M_2 \to M_1, \st hf = \bm{1}_{M_1}$.
    定义
    \begin{align*}
        \Phi: M_2 &\to M_1 \oplus M_3 \\[-5pt]
        m &\mapsto (h(m), g(m))
    \end{align*}

    若$\Phi(m) = 0$，则$g(m) = 0$，$m \in \ker g$. \par
    由于正合列，$m \in \im f$，从而$m = f(m')$. \par
        故$h(m) = h(f(m')) = \bm{1}_{M_1}(m') = m' = 0$. 得$m=0$，故$\Phi$是单射.\par
        下证满射：任取$(m_1, m_3) \in M_1 \oplus M_3$，由于$g$是满射，存在$m' \in M_2$使得$g(m') = m_3$. \par
        取$m = m' + f(m_1 - h(m')) \in M_2$，则有：
        \begin{gather*}
            g(m) = g(m') + g(f(m_1 - h(m'))) = m_3 + 0 = m_3, \\[-5pt]
            h(m) = h(m') + hf(m_1 - h(m')) = h(m') + (m_1 - h(m')) = m_1.
        \end{gather*}
        故$\Phi(m) = (m_1, m_3)$，$\Phi$是满射. 综上所述，$\Phi$是同构.
\end{proof}

既然直和可以通过通过其投影和内射进行分解，若考虑多个模的直积$M_i$，一个自然的问题是：
\begin{align*}
    \Hom\ab(N, \prod_{i\in I} M_i) &\cong \prod_{i\in I1} \Hom(N, M_i) \tag{1}\\
    \Hom\ab(\prod_{i\in I} M_i, M) &\cong \prod_{i\in I} \Hom(M_i, M) \tag{2}
\end{align*}
是否成立？问题的答案是(1)始终成立，但(2)仅在$I$有限时，即直积和直和等价时成立.
为了证明上面两个问题，需要直积的泛性质.

\begin{property}[直积的泛性质]\index{z@直积!f@泛性质}\index{f@泛性质!z@直积}
    设$\pi_j: \prod_{i\in I} M_i \to M_j$为典范投影，则$\forall\, f_j: N \to M_j$，$ \exists\, !\, \overline{f}: N \to \prod_{i\in I} M_i $
    使得下述交换图成立
    \begin{center}
    \begin{tikzcd}[labels={font=\normalsize}]
        N \arrow[r, "f_j"] \arrow[dd, "\exists\, !\, \overline{f}"', dashed] & M_j \\
        \\
        \prod_{i\in I} M_i \arrow[ruu, "\pi_j"']
    \end{tikzcd}
    \end{center}
    即$\pi_j \overline{f} = f_j$.
\end{property}

\begin{proof}
    定义：
    \begin{align*}
        \overline{f}: N &\to \prod_{i\in I} M_i \\[-5pt]
        n &\mapsto \{f_i(n)\}
    \end{align*}
    则$\pi_j \overline{f}(n) = f_j(n)$，满足要求. \par
    若$\overline{f}'$也满足要求，则$\pi_j \overline{f}' = f_j$，$\forall\, j \in I$，则$\overline{f}'(n) = \{f_j(n)\} = \overline{f}(n)$，$\overline{f}' = \overline{f}$.
\end{proof}
利用直积的泛性质可以轻松证明(1)成立.\par
定义
\begin{align*}
    \Phi: \Hom\ab(N, \prod_{i\in I} M_i) &\to \prod_{i\in I} \Hom(N, M_i) \\[-5pt]
    \overline{f} &\mapsto \ab\{\pi_i \circ \overline{f}\}_{i \in I}
\end{align*}

若$\Phi(\overline{f}) = \Phi(\overline{g})$，则对全体 $i \in I$，有 $\pi_i \circ \overline{f} = \pi_i \circ \overline{g}$.由泛性质的唯一性可知 $\overline{f} = \overline{g}$，从而 $\Phi$ 为单射.\par
又因为对于任意的 $\{f_i\} \in \prod_{i\in I} \Hom(N, M_i)$，根据直积泛性质，必定存在唯一的映射 $\overline{f}: N \to \prod_{i\in I} M_i$ 满足 $\pi_i \circ \overline{f} = f_i$，此时恰好有 $\Phi(\overline{f}) = \{f_i\}$，这意味着 $\Phi$ 为满射.\par
综上，$\Phi$ 为同构，即(1)成立.\par

\par


% ---------- 过渡 ----------
下面通过 $\Hom$ 函子的正合性，引入投射模和内射模的定义.
% ---------- 过渡 ----------

设$M, N$是$R-\!\!$模，$\Hom_R(M, N)$为$M \to N$的$R-\!\!$模同态空间. $\forall\, \varphi: A \to B$，可以诱导：
\begin{align*}
    \bar{\varphi}: \Hom_R(M, A) &\to \Hom_R(M, B), & h &\mapsto \varphi h \\
    \bar{f}: \Hom_R(B, N) &\to \Hom_R(A, N), & h &\mapsto h f
\end{align*}

现在，若$\exactseq{0}[]{A}[f]{B}[g]{C}[]{0}$正合，能否推出
\[ \exactseq{0}[]{\Hom_R(M, A)}[\bar{f}]{\Hom_R(M, B)}[\bar{g}]{\Hom_R(M, C)}[]{0} \]
正合，即(1) $\bar{f}$单 (2) $\im \bar{f} = \ker \bar{g}$ (3) $\bar{g}$满.\par
问题的答案是(1)(2)正确，但(3)不满足，即最右侧的$\to 0$应被去掉.
\begin{enumerate}[(1)]
    \item 若$\bar{f}(h) = 0$，则$f h = 0$，由于$f$单，故$h = 0$.
    \item $\im \bar{f} \subseteq \ker \bar{g}$：由于$g f = 0$，故$\bar{g}\bar{f}(h) = g f h = 0$，则$\im \bar{f} \subseteq \ker \bar{g}$.\par
    $\ker \bar{g} \subseteq \im \bar{f}$：
    取$h \in \ker \bar{g}$, $\bar{g}(h) = g h = 0 \Rightarrow \im h \subseteq \ker g = \im f$.\par
    由$f$单，$f': A \to \im f$是同构.
    那么$\bar{f}((f')^{-1} h) = f (f')^{-1} h = h$.
\end{enumerate}

因此实际上我们可以得到以下两个序列的正合性相同：
\begin{proposition}
    \begin{equation*}
    \exactseq{0}[]{A}[f]{B}[g]{C}[]{0}
    \end{equation*}正合当且仅当
    \[ \exactseq{0}[]{\Hom_R(M, A)}[\bar{f}]{\Hom_R(M, B)}[\bar{g}]{\Hom_R(M, C)}[]{0} \]
    正合.
\end{proposition}


同样地，要证(1) $\bar{f}$单，(2) $\im \bar{f} \subseteq \ker \bar{g}$ (3) $\im \bar{f} \supseteq \ker \bar{g}$.
\begin{enumerate}[(1)]
    \item 取$M = \ker f$，有嵌入$\ker f \xrightarrow{\iota} A$.
    $\bar{f}(\iota) = f \iota = 0$，由于单，$\iota = 0$.
    \item 取$M = A$，则$\bar{g}\bar{f}(1) = g f = 0$.
    \item 取$M = \ker g \subseteq B$，有嵌入$\ker g \xrightarrow{\iota} B$.
    $\bar{g}(\iota) = g \iota = 0 \Rightarrow \iota \in \ker \bar{g} = \im \bar{f}$.
    故$\exists\, h \in \Hom(\ker g, A), \st f h = \bar{f}(h) = \iota$.
    因此$f(h(x)) = \iota(x) \in \ker g \Rightarrow \im f \supseteq \ker g$.
\end{enumerate}

\begin{proposition}[]
    \begin{equation*}
    \exactseq{A}[]{B}[]{C}[]{0}
    \end{equation*}正合当且仅当
    \[ 0 \to \Hom_R(C, N) \to \Hom_R(B, N) \to \Hom_R(A, N) \to 0 \]
    正合.
\end{proposition}

称$\Hom(M, -)$为\bluebf{左正合函子}，$\Hom(-, N)$为\bluebf{右正合函子}.\par
为了使上述函子完全正合，可以定义投射模和内射模.


